\documentclass[11pt]{article}
\usepackage[margin=1in]{geometry}
\usepackage{amsmath,amssymb,amsthm,amsfonts,mathtools}
\usepackage{algorithm,algpseudocode}
\usepackage{booktabs}
\usepackage{hyperref}
\usepackage{xcolor}
\usepackage{graphicx}

% ── Theorem environments ──────────────────────────────────────────────────────
\theoremstyle{plain}
\newtheorem{definition}{Definition}
\newtheorem{lemma}{Lemma}
\newtheorem{proposition}{Proposition}
\newtheorem{theorem}{Theorem}
\newtheorem{corollary}{Corollary}
\theoremstyle{remark}
\newtheorem{remark}{Remark}

% ── Notation shorthands ───────────────────────────────────────────────────────
\newcommand{\supp}{\operatorname{supp}}
\newcommand{\prods}{\operatorname{prod}}
\newcommand{\cl}{\operatorname{cl}}
\newcommand{\MinBas}{\operatorname{MinBas}}
\newcommand{\req}{\operatorname{req}}
\newcommand{\supl}{\operatorname{supl}}
\newcommand{\minprod}{\operatorname{minprod}}
\newcommand{\mincons}{\operatorname{mincons}}
\newcommand{\E}{\mathcal{E}}
\newcommand{\R}{\mathcal{R}}
\newcommand{\M}{M}
\newcommand{\N}{\mathbb{N}}
\renewcommand{\S}{S}
\newcommand{\js}{\vee}
\newcommand{\SO}{\mathit{SO}}
\newcommand{\SOgen}{\mathit{SO}^{\mathrm{gen}}}
\newcommand{\ord}{\operatorname{ord}}
\newcommand{\fcomp}[1]{\rightleftharpoons_{#1}}
\newcommand{\notes}[1]{\textcolor{blue}{[#1]}}

% ─────────────────────────────────────────────────────────────────────────────
\begin{document}

\title{Bitset Algorithms for the Generative Structure of Chemical
Organizations: From Elementary Reaction Closures to All Persistent Modules}

\author{%
  Tom\'as Veloz$^{1,\ast}$ \and Paolo Bassi$^{2}$\\[4pt]
  \small $^{1}$Templeton World Charity Foundation Postdoctoral Fellow,
  CLEA, Vrije Universiteit Brussel, Belgium\\
  \small $^{2}$\ldots\\
  \small $^{\ast}$\texttt{tveloz@gmail.com}
}
\date{}

\maketitle

\begin{abstract}
\noindent\textbf{Motivation:} Chemical Organization Theory (COT) identifies
the closed, self-maintaining subnetworks of a reaction network --
\emph{persistent modules} -- that constrain its long-run dynamical
behaviour. Deciding whether a network contains one is NP-hard, and the
established bottom-up algorithms enumerate the lattice of closed sets,
most of which are generatively irrelevant, so exact computation has been
confined to networks of a few dozen reactions while genome-scale metabolic
reconstructions reach thousands. A companion theoretical study shows that
every persistent module is generated from \emph{elementary reaction
closures} (ERCs) by two irreducible relations, fundamental synergy and
fundamental complementarity, but stops at an existence theorem.

\noindent\textbf{Results:} We give and implement an algorithm that
computes, exactly, all elementary persistent modules (EPMs) and all
elementary-supported persistent modules (ESPMs) of a reaction network from
its ERCs. The search is carried out entirely in ERC-index space: each
partial module is three incrementally updated bitmasks (species present,
species produced, species still externally required), so no reaction scan
or closure recomputation is ever needed after the ERC hierarchy, the
fundamental-synergy hypergraph and the fundamental-complementarity graph
have been built once. A module grows by three moves -- attaching a minimal
producer of an open requirement, attaching a synergy partner, or, once the
module is already self-maintaining, attaching a minimal consumer of a
species it already produces or lifting to a hierarchy ancestor of one of
its constituents -- and the algorithm deliberately searches only this
\emph{generated core}: unions of already-persistent modules that share no
synergy or complementarity (\emph{latent joins}) are excluded from the
search and reconstructed in a single closure computation on demand. We
validate correctness by exhaustive comparison against a brute-force
oracle on hand-built networks, a fully worked network exhibiting every
mechanism including a sterile ERC, and a sample of real BioModels networks
up to the size at which brute force remains tractable, finding exact
agreement throughout. On a genome-scale metabolic reconstruction (825
species, 1296 reactions, 712 ERCs) the algorithm finds 72 EPMs and 1124
ESPMs across 17 orders in under 20 seconds on one desktop core -- a scale
at which brute-force enumeration is entirely infeasible.

\noindent\textbf{Availability:} Implemented in the open-source
\texttt{cot\_gen} library; \url{https://github.com/tveloz/pyCOT}.

\noindent\textbf{Contact:} \texttt{tveloz@gmail.com}
\end{abstract}

\vspace{4pt}
\noindent\textbf{Keywords:} Chemical Organization Theory, reaction
networks, elementary reaction closures, synergy, complementarity,
persistent modules, bitset algorithms.

% =============================================================================
\section{Introduction}\label{sec:intro}
% =============================================================================

Chemical Organization Theory (COT) studies the long-run behaviour of a
reaction network by identifying its closed, self-maintaining subsets, or
\emph{organizations}~\cite{dittrich2007chemical}. A companion theoretical
study~\cite{BMBcompanion} isolates the minimal building blocks from which
every such subset is assembled -- \emph{elementary reaction closures}
(ERCs) -- organises them into a containment hierarchy, and identifies two
relations between ERCs, \emph{synergy} (a combination fires a reaction
neither part fires alone) and \emph{complementarity} (one part produces a
species the other requires), whose most restrictive, \emph{fundamental}
forms are proven sufficient to generate every persistent module. That
study proves existence but gives neither an algorithm nor an
implementation: it stops at brute-force statistics over $438$ biological
networks.

This paper supplies the missing computational layer. We describe, and have
implemented and validated, the \texttt{cot\_gen} library, which turns the
companion study's generative theorems into a concrete pipeline: it
discovers ERCs and their hierarchy, computes basic/maximal/fundamental
synergies and basic/pure/fundamental complementarities, identifies the
primitive ERCs and their synergy-generated reach, and -- the central
contribution of this paper -- computes \emph{all} elementary persistent
modules (EPMs) and elementary-supported persistent modules (ESPMs) of a
reaction network, of any order, by a bounded-branching search in ERC-index
space rather than by enumerating the lattice of closed sets. Every stage
is cross-validated against an independent brute-force oracle, on
hand-built networks, on a fully worked example, and on a sample of real
BioModels networks up to the size at which the oracle itself remains
tractable.

\paragraph{Structure.} Section~\ref{sec:background} reviews the COT
definitions this paper builds on. Sections~\ref{sec:stage1}--\ref{sec:stage5}
describe ERC discovery, the containment hierarchy, synergy, complementarity,
and the primitive-ERC generative basis. Section~\ref{sec:grammar} gives the
central contribution: the ERC-index-space representation and the two-mode
traversal that computes all EPMs and ESPMs while searching only the
generated core. Section~\ref{sec:implementation} describes the bitset
implementation. Section~\ref{sec:oracle} describes the oracle-first
validation protocol and its results, including a fully worked example.
Section~\ref{sec:experiments} reports a representative genome-scale run.
Section~\ref{sec:conclusion} concludes and lists the algorithmic
extensions this implementation makes possible.

% =============================================================================
\section{Background}\label{sec:background}
% =============================================================================

\subsection{Reaction Networks and the $E_0$ Quotient}

A \emph{reaction network} (RN) is a tuple $(\M, \R)$ where
$\M = \{s_1,\ldots,s_n\}$ is a finite set of \emph{species} and
$\R = \{r_1,\ldots,r_m\}$ is a finite set of \emph{reactions}.
Each reaction $r$ has a \emph{support} $\supp(r) \subseteq \M$
(species consumed) and a \emph{product set} $\prods(r) \subseteq \M$
(species produced).

\begin{definition}[$E_0$ set]
The \emph{inflow set} $E_0 \subseteq \M$ is the closure of all species
produced by reactions with empty support (\emph{inflow reactions}):
$E_0 = \cl(\bigcup_{r:\supp(r)=\emptyset}\prods(r))$.
\end{definition}

All subsequent computation uses the \emph{quotiented} support and
products, $\supp_q(r) = \supp(r) \setminus E_0$ and
$\prods_q(r) = \prods(r) \setminus E_0$; reactions with
$\supp_q(r) = \emptyset$ play no further role in ERC computation, since
$E_0$ already sits inside every persistent module and is suppressed
below.

\begin{definition}[Closure]
The \emph{closure} $\cl(X)$ of a species set $X \subseteq \M$ is the
smallest $Y \supseteq X$ such that
$\supp_q(r) \subseteq Y \Rightarrow \prods_q(r) \subseteq Y$ for all $r$.
\end{definition}

\subsection{Elementary Reaction Closures}

\begin{definition}[ERC~\cite{BMBcompanion}]
For each reaction $r$ with $\supp_q(r) \neq \emptyset$, its
\emph{ERC mask} is $E_r = \cl(\supp_q(r))$.
The ERC set is $\E = \{E_r : r \in \R,\, \supp_q(r) \neq \emptyset\}$.
Distinct reactions with the same closure are grouped: $\E$ is the set of
\emph{distinct} closures.
\end{definition}

\begin{definition}[Requirement and production masks]
For $E \in \E$ let $\R_E = \{r : \supp_q(r) \subseteq E\}$ be the
reactions active within $E$. Then
\[
  \prods(E) = \bigcup_{r \in \R_E} \prods_q(r), \qquad
  \req(E) = \Bigl(\bigcup_{r \in \R_E} \supp_q(r)\Bigr) \setminus \prods(E).
\]
$E$ is \emph{persistent} (a P-ERC) iff $\req(E) = \emptyset$.
\end{definition}

ERCs are partially ordered by inclusion; the Hasse diagram of direct
containments is a forest~\cite{BMBcompanion}.

\subsection{Semi-Organizations and Persistent Modules}

\begin{definition}[Semi-organization (SO)~\cite{BMBcompanion}]
A species set $X \subseteq \M$ is a \emph{semi-organization} iff it is
(i)~\emph{closed}: $\prods(\R_X) \subseteq X$;
(ii)~\emph{semi-self-maintaining} (SSM): $\req(X) = \emptyset$;
(iii)~\emph{reactive}: $\R_X \neq \emptyset$;
(iv)~\emph{connected}: two species of $X$ are connected iff they co-occur
in the support or product of some active reaction, and this relation
spans all of $X$.
We write $\SO$ for the set of all semi-organizations of a network, our
computational target.
\end{definition}

\begin{definition}[EPM and ESPM order~\cite{BMBcompanion}]
\label{def:order}
$X \in \SO$ is an \emph{elementary persistent module} (EPM) if its only
proper sub-SO is $E_\emptyset$ (the empty set, once $E_0$ is quotiented
out); write $\ord(X)=0$. Otherwise $X$ is an \emph{elementary-supported
persistent module} (ESPM) of order
$\ord(X)=1+\max\{\ord(Y): Y \in \SO,\ \emptyset \subsetneq Y \subsetneq X\}$.
\end{definition}

\subsection{Synergy and Complementarity}

Every reactive closed set is a join of ERCs, so ERCs generate $\SO$
\cite{BMBcompanion}. Two ERCs are \emph{synergetic} if their join fires a
reaction neither fires alone; and \emph{complementary} if one supplies a
requirement of the other, measured by
$\supl(E \rightharpoonup E') = \prods(E) \cap \req(E')$. The
\emph{fundamental} forms of each relation minimise the participants (a
fundamental synergy uses the smallest reactant ERCs producing the largest
target; a fundamental complementarity carries a species from a
$\subseteq$-minimal producer to a $\subseteq$-minimal consumer -- exact
definitions in Sections~\ref{sec:stage3} and~\ref{sec:stage4}). The
companion study proves that every $X \in \SO$ admits a \emph{fundamental
generator}: an irreducible sequence of ERC attachments from
$E_\emptyset$, each of which is either a fundamental synergy or a
fundamental complementarity, and that this suffices to reach every
persistent module without ever forming a reaction-inert combination of
unrelated closed sets~\cite{BMBcompanion}. Sections~\ref{sec:stage1}
through~\ref{sec:stage5} give the bitset algorithms that compute the ERC
hierarchy and these fundamental relations; Section~\ref{sec:grammar} gives
the algorithm that walks them to enumerate $\SO$ itself.

% =============================================================================
\section{ERC Discovery}\label{sec:stage1}
% =============================================================================

\begin{algorithm}[t]
\caption{ComputeERCs($\R$, $E_0$)}\label{alg:ercs}
\begin{algorithmic}[1]
\State $\text{groups} \leftarrow \{\}$ \Comment{mask $\mapsto$ list of reaction indices}
\For{each $r \in \R$ with $\supp_q(r) \neq \emptyset$}
  \State $E \leftarrow \cl(\supp_q(r))$
  \State $\text{groups}[E].\text{append}(r)$
\EndFor
\State $\text{ercs} \leftarrow []$
\For{each $(E, \text{rxns})$ in groups}
  \State $\text{bases} \leftarrow \MinBas(\{\supp_q(r) : r \in \text{rxns}\})$
  \State Compute $\req(E)$, $\prods(E)$ from $\R_E$
  \State $\text{ercs.append}(\text{ERCData}(E, \text{rxns}, \text{bases}, \req(E), \prods(E)))$
\EndFor
\State \Return ercs sorted by $|E|$
\end{algorithmic}
\end{algorithm}

Closure is computed by a Horn/Dowling--Gallier-style
propagation~\cite{dowling1984linear}: a reactant counter per reaction is
decremented as species are added, and a reaction fires the instant its
counter reaches zero, avoiding the naive fixpoint's repeated rescans.
Each $\MinBas(E)$ is the $\subseteq$-minimal antichain of reaction
supports whose closure is $E$; a species set ignites $E$ iff it contains
some element of $\MinBas(E)$ (Lemma~\ref{lem:ignition} below), which is
what makes bases the right domain for detecting synergy in
Section~\ref{sec:stage3}.

\paragraph{Complexity.} Closure via Horn propagation costs
$O(\sum_r|\supp_q(r)|)$ amortised per call; $O(|\R|)$ closure calls give
$O(|\R|\cdot\sum_r|\supp_q(r)|)$ total. The $\MinBas$ antichain filter
costs $O(k^2)$ per group of $k$ same-closure reactions.

% =============================================================================
\section{Containment Hierarchy}\label{sec:stagehier}
% =============================================================================

\begin{definition}[ERC order]
$E_i \leq E_j \iff E_i \subseteq E_j$ as species bitmasks.
\end{definition}

The Hasse diagram stores only direct-cover edges (no intermediate ERC
between them); the algorithm first computes all containment pairs via the
bitset predicate $(m_i\,\&\,m_j) = m_i$, then removes transitively implied
edges, and finally derives, for every ERC, its full ancestor and
descendant sets by transitive closure -- these are the sets the traversal
of Section~\ref{sec:grammar} queries directly (\texttt{parents},
\texttt{ancestors}, \texttt{descendants}).

\paragraph{Complexity.} $O(|\E|^2)$ for the pairwise containment scan and
Hasse reduction; $O(|\E|^2)$ for the ancestor/descendant closure.

% =============================================================================
\section{Synergy}\label{sec:stage3}
% =============================================================================

\subsection{Definitions~\cite{BMBcompanion}}

\begin{definition}[Basic synergy]
$(E_i, E_j) \rightsquigarrow E_k$ is a \emph{basic synergy} iff
$E_i, E_j, E_k$ are pairwise incomparable and
$\exists\, b \in \MinBas(E_k)$ with
$b \subseteq (E_i \cup E_j)$, $b \not\subseteq E_i$, $b \not\subseteq E_j$:
the joint species set ignites a basis of $E_k$ that neither part ignites
alone.
\end{definition}

\begin{definition}[Maximal synergy]
$(E_i, E_j) \rightsquigarrow^{\max} E_k$ is \emph{maximal} iff it is basic
and no $E_{k'} \supsetneq E_k$ is also a basic synergy target of the pair
$(E_i, E_j)$.
\end{definition}

\begin{definition}[Fundamental synergy]
\label{def:fund_syn}
A maximal synergy $(E_i, E_j) \rightsquigarrow^{\max} E_k$ is
\emph{fundamental} (written $E_i+E_j\twoheadrightarrow E_k$) iff both
reactant slots are irreducible: no $E_{i'} \subsetneq E_i$ gives
$(E_{i'}, E_j) \rightsquigarrow^{\max} E_k$, and no $E_{j'} \subsetneq E_j$
gives $(E_i, E_{j'}) \rightsquigarrow^{\max} E_k$.
\end{definition}

\subsection{Pair-First Algorithm}

\begin{algorithm}[t]
\caption{ComputeFundamentalSynergies-PairFirst($\E$, hier)}\label{alg:syn}
\begin{algorithmic}[1]
\For{all incomparable pairs $(i, j)$}
  \For{each $k$ not a descendant of $i$ or $j$}
    \If{$\exists\, b \in \MinBas(k):\; b \subseteq m_i|m_j,\; b\not\subseteq m_i,\; b\not\subseteq m_j$}
      \State $\text{basic}[i,j].\text{add}(k)$
    \EndIf
  \EndFor
\EndFor
\State $\text{maximal} \leftarrow \{(i,j,k) \in \text{basic} :
  \nexists k'\ \text{with}\ m_k \subsetneq m_{k'} \in \text{basic}[i,j]\}$
\For{each $(i,j,k) \in \text{maximal}$}
  \If{$\nexists i':\; m_{i'} \subsetneq m_i \land (i',j,k) \in \text{maximal}$
      \textbf{and} $\nexists j':\; m_{j'} \subsetneq m_j \land (i,j',k) \in \text{maximal}$}
    \State emit fundamental$(i,j,k)$
  \EndIf
\EndFor
\end{algorithmic}
\end{algorithm}

\paragraph{Complexity.} $O(|\E|^3 \cdot B)$ where
$B = \max_k|\MinBas(E_k)|$; used as the reference implementation and for
cross-validating the basis-first algorithm below.

\subsection{Basis-First Algorithm (default)}

The pair-first scan checks every incomparable pair against every target's
bases, $O(|\E|^3\!\cdot\!B)$ overall, even though on sparse networks most
ERC pairs share no species with any given basis. Basis-first inverts the
loop order: for each target $k$ and each basis $b\in\MinBas(k)$, only the
ERCs that cover \emph{some but not all} of $b$ are relevant reactant
candidates -- an ERC covering all of $b$ ignites $k$ by itself and is not
a synergy source, and one covering none of $b$ contributes nothing to it.

\begin{algorithm}[t]
\caption{ComputeFundamentalSynergies-BasisFirst($\E$, hier)}\label{alg:syn-basis}
\begin{algorithmic}[1]
\For{each target $k \in \E$}
  \For{each basis $b \in \MinBas(k)$}
    \State $\text{relevant} \leftarrow \{i : \emptyset \subsetneq (b\,\&\,m_i) \subsetneq b\}$
    \For{each incomparable pair $(i,j)$ in relevant with $(b\,\&\,m_i)\,|\,(b\,\&\,m_j)=b$}
      \State $\text{basic}[i,j].\text{add}(k)$
    \EndFor
  \EndFor
\EndFor
\State apply the same maximal and per-slot fundamentality filters as Algorithm~\ref{alg:syn}
\end{algorithmic}
\end{algorithm}

\begin{proposition}[Equivalence]
Algorithms~\ref{alg:syn} and~\ref{alg:syn-basis} report the same basic
synergy set. For $(i,j)\to k$: let $b_i = b\,\&\,m_i$, $b_j = b\,\&\,m_j$.
$(i,j)\rightsquigarrow k$ via $b$ in the pair-first sense iff
$b_i\neq\emptyset$, $b_j\neq\emptyset$, $b_i\neq b$, $b_j\neq b$, and
$b_i\,|\,b_j = b$ -- exactly the basis-first membership test, by De
Morgan's laws on the joint-cover predicate.
\end{proposition}

The maximal and fundamental filters are reused unchanged from
Algorithm~\ref{alg:syn}. Basis-first is the algorithm used by default
throughout this paper and by \texttt{cot\_gen}'s pipeline script; the
pair-first algorithm is retained purely as an independent cross-check
(Section~\ref{sec:oracle}).

\paragraph{Complexity.} Bounded by $O(|\E|\cdot B \cdot \max_b|\text{relevant}(b)|^2)$,
which is markedly smaller than $O(|\E|^3\!\cdot\!B)$ whenever bases are
locally concentrated -- the typical case for metabolic and signalling
networks (Section~\ref{sec:experiments}).

% =============================================================================
\section{Complementarity}\label{sec:stage4}
% =============================================================================

\subsection{Definitions~\cite{BMBcompanion}}

\begin{definition}[Supply]
The \emph{supply} from $E$ to $E'$ is
$\supl(E \rightharpoonup E') = \prods(E) \cap \req(E')$.
In bitsets: $\mathtt{prod\_mask}_E \;\&\; \mathtt{req\_mask}_{E'}$.
\end{definition}

\begin{definition}[Basic complementarity]
$(E, E')$ are \emph{basic complementary} iff
$\supl(E \rightharpoonup E') \cup \supl(E' \rightharpoonup E) \neq \emptyset$:
at least one ERC produces a species the other requires. Only incomparable
pairs are considered, since for $E\subsetneq E'$, $\prods(E)\subseteq\prods(E')$
gives $\supl(E\rightharpoonup E')=\emptyset$ trivially.
\end{definition}

\begin{definition}[Pure complementarity]
$(E, E')$ are \emph{purely complementary} iff they are basic complementary
and \emph{not} synergetic (no basic synergy involves the pair). This
isolates requirement reduction attributable directly to ERC-pair supply
from requirement reduction that is really an artefact of a jointly
activated reaction -- see Remark~\ref{rem:gold2} below.
\end{definition}

\begin{definition}[Fundamental complementarity]
\label{def:fund_comp}
For species $s \in \supl(E \rightharpoonup E')$, the supply is
\emph{fundamental} iff $E \in \minprod(s)$ and $E' \in \mincons(s)$, where
$\minprod(s)=\min_\subseteq\{E\in\E: s\in\prods(E)\}$ and
$\mincons(s)=\min_\subseteq\{E\in\E: s\in\req(E)\}$ are the
$\subseteq$-minimal producers and consumers of $s$. We write
$E\fcomp{s}E'$.
\end{definition}

\begin{remark}[Complementarity does not certify closed-set-level requirement
reduction]
\label{rem:gold2}
Joining two persistent-module constituents can strictly reduce their
combined requirement, $\req(X\js Y)\subsetneq\req(X)\cup\req(Y)$, purely
because the join fires a new (synergetic) reaction whose product closes a
requirement that neither part supplied on its own; the pair itself then
carries \emph{zero} basic complementarity under Definition~4, even though
the module as a whole becomes more self-sufficient. This is precisely why
pure complementarity and fundamental synergy are tracked as separate,
non-overlapping relations rather than folded into a single ``requirement
shrank'' predicate, and it is exercised directly by our test suite
(Section~\ref{sec:oracle}).
\end{remark}

\subsection{Algorithms}

\begin{algorithm}[t]
\caption{ComputeBasicComplementarity($\E$, hier)}\label{alg:basic_comp}
\begin{algorithmic}[1]
\For{all incomparable pairs $(i, j)$}
  \State $\text{fwd} \leftarrow p_i \;\&\; r_j$; \quad
         $\text{bwd} \leftarrow p_j \;\&\; r_i$
  \If{$\text{fwd} \neq 0$ or $\text{bwd} \neq 0$}
    \State emit BasicComp$(i, j, \text{fwd}, \text{bwd},\ \text{is\_pure})$
  \EndIf
\EndFor
\end{algorithmic}
\end{algorithm}

\begin{algorithm}[t]
\caption{ComputeFundamentalComplementarity($\E$)}\label{alg:fund_comp}
\begin{algorithmic}[1]
\For{each species bit $s$}
  \State $\text{prod\_s}, \text{req\_s} \leftarrow$ producers, requirers of $s$
  \If{$\text{prod\_s} = \emptyset$ or $\text{req\_s} = \emptyset$}
    \State \textbf{continue}
  \EndIf
  \State $\minprod(s), \mincons(s) \leftarrow$ $\subseteq$-minimal antichains of prod\_s, req\_s
  \For{$\pi \in \minprod(s)$, $\gamma \in \mincons(s)$, $\pi \neq \gamma$}
    \State emit FundComp$(\pi, \gamma, s)$
    \Comment{indexed by both producer $\pi$ and consumer $\gamma$}
  \EndFor
\EndFor
\end{algorithmic}
\end{algorithm}

Algorithm~\ref{alg:fund_comp} builds two inverted indices over the emitted
\texttt{FundComp} triples -- \texttt{comp\_by\_species}$[s]$, the minimal
producers of $s$, and \texttt{comp\_consumers\_by\_species}$[s]$, the
minimal consumers of $s$ -- because the traversal of
Section~\ref{sec:grammar} needs both directions: producers to close an
open requirement, consumers to discharge productive novelty once a module
is already self-maintaining.

\paragraph{Complexity.} Algorithm~\ref{alg:basic_comp}: $O(|\E|^2)$.
Algorithm~\ref{alg:fund_comp}: $O(|\M| \cdot |\E|^2)$, dominated by the
$\minprod/\mincons$ antichain filters.

% =============================================================================
\section{Primitive ERCs and the Generative Basis}\label{sec:stage5}
% =============================================================================

\begin{definition}[Primitive ERC]
$E \in \E$ is \emph{primitive} iff it is not the target of any fundamental
synergy.
\end{definition}

\begin{definition}[Synergy reachability]
The \emph{basis reach} $\mathcal{B}$ is the smallest superset of the
primitive ERCs $\mathcal{P}$ closed under fundamental synergy: if
$E_i,E_j\in\mathcal{B}$ and $E_i+E_j\twoheadrightarrow E_k$, then
$E_k\in\mathcal{B}$.
\end{definition}

\begin{algorithm}[t]
\caption{ComputeGenerators($\E$, $\mathcal{F}$)}\label{alg:gen}
\begin{algorithmic}[1]
\State $\mathcal{P} \leftarrow \{E \in \E : E \notin \{E_k : (i,j,k) \in \mathcal{F}\}\}$
\State $\mathcal{B} \leftarrow \mathcal{P}$
\While{$\mathcal{B}$ changes}
  \For{each $(i,j,k) \in \mathcal{F}$}
    \If{$E_i, E_j \in \mathcal{B}$}
      \State $\mathcal{B} \leftarrow \mathcal{B} \cup \{E_k\}$
    \EndIf
  \EndFor
\EndWhile
\State \Return $(\mathcal{P},\; \mathcal{B},\; |\mathcal{B}|/|\E|)$
\end{algorithmic}
\end{algorithm}

This synergy-only reachability closure identifies how much of the ERC
poset is explained purely by combining primitives via fundamental synergy;
it is a diagnostic on $\E$ (Section~\ref{sec:experiments} reports
coverage on a genome-scale network) and is distinct from, and computed
before, the full module search of Section~\ref{sec:grammar}, which also
draws on complementarity and hierarchy structure.

\paragraph{Complexity.} $O(|\E| \cdot |\mathcal{F}|)$.

% =============================================================================
\section{Computing All Persistent Modules}\label{sec:grammar}
% =============================================================================

This section gives the paper's central contribution: an algorithm that
computes \emph{every} EPM and ESPM of a reaction network -- not just
single-ERC or pairwise combinations -- directly from the structures of
Sections~\ref{sec:stage1}--\ref{sec:stage5}, without ever computing a
reaction-level closure during the search itself.

\subsection{Modules as ERC-index-space states}

\begin{lemma}[Ignition]
\label{lem:ignition}
For closed $X$ and $E\in\E$, $E\subseteq X$ iff some $B\in\MinBas(E)$ has
$B\subseteq X$.
\end{lemma}
\begin{proof}
($\Rightarrow$) any $B\subseteq E\subseteq X$. ($\Leftarrow$) $B\subseteq X$
closed gives $E=\cl(B)\subseteq X$.
\end{proof}

A partial module under construction is represented not as a species set
but as a set $S\subseteq\{1,\dots,|\E|\}$ of ERC indices, together with
three incrementally maintained bitmasks:
\[
  \mathrm{sp}(S)=\bigcup_{i\in S}E_i,\qquad
  \mathrm{prod}(S)=\bigcup_{i\in S}\prods(E_i),\qquad
  \mathrm{req}(S)=\Bigl(\bigcup_{i\in S}\req(E_i)\Bigr)\setminus\mathrm{prod}(S).
\]
These formulae are exact -- not approximations -- because $\req(E_i)$
already incorporates every reaction internal to $E_i$, including those of
every sub-ERC $E_j\subsetneq E_i$ (since $\R_{E_j}\subseteq\R_{E_i}$ by
containment); the only cross-boundary reactions a combination can newly
activate are exactly the ones captured by the fundamental-synergy
hypergraph. $S$ is self-maintaining iff $\mathrm{req}(S)=\emptyset$; no
separate SSM check or closure computation is ever needed.

\begin{lemma}[Synergy ignition is a species-subset test, not a membership
test]
\label{lem:synignite}
When ERC set $S$ grows to $S\cup\{j\}$, a fundamental synergy
$E_a+E_b\twoheadrightarrow E_k$ fires (i.e.\ $k$ is implied) iff
$E_a\subseteq\mathrm{sp}(S\cup\{j\})$ and $E_b\subseteq\mathrm{sp}(S\cup\{j\})$
-- \emph{not} iff $a,b\in S\cup\{j\}$. These differ whenever a bigger ERC
is present that species-wise dominates a minimal reactant ERC without that
reactant itself ever being an explicit member of $S$: e.g.\ after a
vertical lift (Section~\ref{subsec:moves}) injects an ancestor $E$ of some
$E_i\in S$, every descendant $d$ of $E$ satisfies
$E_d\subseteq E\subseteq\mathrm{sp}(S)$ and so is ignited, even though
$d\notin S$.
\end{lemma}

Lemma~\ref{lem:synignite} is not a stylistic nuance: an implementation
that tests $a,b\in S$ instead of the species-subset condition silently
produces states whose tracked $(\mathrm{sp},\mathrm{req},\mathrm{prod})$
disagree with the true reaction-level closure of $\mathrm{sp}(S)$ -- we
found and fixed exactly this defect during validation
(Section~\ref{sec:oracle}), and it is the reason
Algorithm~\ref{alg:erc-syn-close} below seeds its propagation queue not
only from the literally added indices but from their full hierarchy
descendant sets.

\begin{algorithm}[t]
\caption{ExtendState($S$, $\mathrm{sp}$, $\mathrm{req}$, $\mathrm{prod}$, $j$)}
\label{alg:erc-syn-close}
\begin{algorithmic}[1]
\State $\text{implied} \leftarrow \{j\}$; $\text{sp}' \leftarrow \mathrm{sp}\,|\,m_j$
\State $\text{queue} \leftarrow \{j\}\cup\mathrm{descendants}(j)$
\While{queue nonempty}
  \State pop $\ell$ from queue
  \For{each $(a,b,k)$ with $\ell\in\{a,b\}$, $k\notin\text{implied}$}
    \State let $o$ be the other reactant slot ($a$ or $b$)
    \If{$m_o\,\&\,\text{sp}' = m_o$} \Comment{species-subset ignition test}
      \State $\text{implied}\leftarrow\text{implied}\cup\{k\}$
      \State $\text{sp}'\leftarrow\text{sp}'\,|\,m_k$
      \State push $k$ and $\mathrm{descendants}(k)$ onto queue
    \EndIf
  \EndFor
\EndWhile
\State $\text{add\_req},\text{add\_prod}\leftarrow \bigcup_{i\in\text{implied}}\req(E_i),\ \bigcup_{i\in\text{implied}}\prods(E_i)$
\State $\text{prod}'\leftarrow\mathrm{prod}\,|\,\text{add\_prod}$
\State \Return $(S\cup\text{implied},\ \text{sp}',\ (\mathrm{req}\,|\,\text{add\_req})\,\&\,\sim\text{prod}',\ \text{prod}')$
\end{algorithmic}
\end{algorithm}

Deduplication across the whole search is a single global set of visited
$\mathrm{sp}$ bitmasks: two different ERC-index paths that reach the same
species set are recognised as the same module the instant the second one
is popped, with no separate ``signature'' object needed beyond the species
bitmask itself.

\subsection{Two modes, three extension moves}
\label{subsec:moves}

The search proceeds in two modes, mirroring the productive-novelty
distinction of the companion theory: an open requirement can only be
closed by attaching a \emph{producer}, whereas a module that is already
self-maintaining can only be grown further by attaching a
\emph{consumer} of something it already makes, by lifting to a bigger
ERC, or by synergy.

\begin{description}
\item[Mode 1 (requirement-driven).] While $\mathrm{req}(S)\neq\emptyset$,
for each still-required species $s$, attach every
$\pi\in\texttt{comp\_by\_species}[s]$ (a $\subseteq$-minimal producer of
$s$); and for each $E_i\in S$, for each fundamental-synergy partner $j$ of
$i$ whose target's production intersects $\mathrm{req}(S)$, attach $j$.
Mode 1 halts the branch the instant $\mathrm{req}(S)=\emptyset$
(a candidate SSM) or no attachment is possible (a dead end).
\item[Mode 2 (growth past self-maintenance).] Given an already
self-maintaining $S$, three moves seed further growth, each resumed
through Mode 1 in case it reopens a requirement:
\begin{enumerate}
\item \emph{Synergy}: for $E_i\in S$, attach a fundamental-synergy
partner $j\notin S$.
\item \emph{Vertical lift}: for $E_i\in S$, attach a direct hierarchy
parent $E$ of $E_i$ ($E\supsetneq E_i$). Since $\req(E)$ already
incorporates $E_i$'s reactions, this is exactly another call to
Algorithm~\ref{alg:erc-syn-close} -- no separate ``absorb'' step is
needed.
\item \emph{Consumer complementarity}: for each species $s$ already
produced by $S$ ($s\in\mathrm{prod}(S)$), attach a $\subseteq$-minimal
consumer $\gamma\in\texttt{comp\_consumers\_by\_species}[s]$,
$\gamma\notin S$. This is the mirror image of Mode 1's producer
attachment, using the same fundamental-complementarity relation from its
other side.
\end{enumerate}
\end{description}

\begin{proposition}[Producer factorisation~\cite{BMBcompanion}]
Every requirement-closing attachment $S\to S\cup\{i\}$ with
$s\in\prods(E_i)\cap\req(S)$ factors as a fundamental horizontal core
followed by vertical lifts: producers of $s$ form an up-set, so $E_i$
dominates some $\hat E\in\minprod(s)$; attaching $\hat E$ is a fundamental
complementary extension, and $E_i$ is recovered by climbing
$\hat E\to\cdots\to E_i$. When $\hat E$ is already an ancestor of a
constituent, the horizontal step is vacuous and the attachment is purely
vertical.
\end{proposition}

This is exactly why Mode 1 restricts producer attachment to $\subseteq$-minimal
producers: any larger, non-minimal producer that could also close the
requirement is reachable afterwards, if at all needed, by vertical lift
from the minimal one -- confining the branching factor to the fundamental
relations rather than to $|\E|$ or $\binom{|\E|}{k}$.

\begin{algorithm}[t]
\caption{ComputeAllPersistentModules($\E$, hier, $\mathcal{F}_{\mathrm{syn}}$, $\mathcal{F}_{\mathrm{comp}}$)}
\label{alg:main}
\begin{algorithmic}[1]
\State \textbf{// Single-ERC EPMs (no search needed)}
\State $\mathrm{EPM}_1 \leftarrow \{E_i : \req(E_i)=\emptyset,\ \nexists\,E_j\subsetneq E_i,\ \req(E_j)=\emptyset\}$
\State \textbf{// Mode 1: one canonical-order seed per non-persistent ERC}
\For{each non-persistent $E_i$}
  \State seed a DFS state $\{i\}$ with $\mathrm{min\_ext}=i{+}1$
  \Comment{only extensions $\geq i{+}1$ allowed -- see Remark~\ref{rem:canon}}
\EndFor
\State run Mode 1 from every seed until every branch is an SSM or a dead end; collect $\mathrm{SSM}_1$
\State \textbf{// Order assignment (post-hoc, by species-mask containment)}
\State sort $\mathrm{EPM}_1\cup\mathrm{SSM}_1$ by popcount; assign
       $\ord(X)=1+\max\{\ord(Y): Y\subsetneq X,\ Y \text{ already ordered}\}$, else $0$
\State \textbf{// Mode 2: BFS by order}
\For{$k=1,2,\dots$ until no new module}
  \For{each $X_p$ with $\ord(X_p)=k-1$}
    \State seed Mode 2 extensions (synergy, vertical lift, consumer complementarity)
    \State resume Mode 1 from each seed; new SSMs get order via the same containment rule
  \EndFor
\EndFor
\State \Return all discovered SOs, grouped by order
\end{algorithmic}
\end{algorithm}

\begin{remark}[Canonical ordering]
\label{rem:canon}
A DFS seeded from every non-persistent ERC would rediscover every
multi-ERC SSM once per permutation of its constituents. Each seed state
carries a $\mathrm{min\_ext}$ gate: an explicit attachment must have index
$\geq\mathrm{min\_ext}$, and any synergy-implied attachment with index
below the seed's own smallest constituent prunes the branch (its
canonical path is from that smaller seed instead). This confines every
$k!$-fold permutation redundancy to a single canonical discovery, at the
cost of no completeness: every SSM is still found, from its smallest-index
constituent.
\end{remark}

\subsection{The generated core, and latent joins on demand}

\begin{definition}[Latent join]
$X\in\SO$ is \emph{latent} if some irreducible generator of $X$ consists
entirely of steps that neither fire a new reaction nor strictly reduce a
requirement -- equivalently, $X$ is the join of strictly smaller SOs
sharing no synergy or complementarity between them. Write $\SOgen$ for the
non-latent (\emph{generated}) modules.
\end{definition}

\begin{proposition}[Latent--generated factorisation~\cite{BMBcompanion}]
$\SO=\langle\SOgen\rangle_\js$: every SO is a join of generated modules,
and because $\req(X\js Y)\subseteq\req(X)\cup\req(Y)$ always, joining two
SSM modules is automatically SSM -- so nothing about a latent join needs
to be searched for; it is a single closure computation away from any two
of its generated pieces.
\end{proposition}

Algorithm~\ref{alg:main} therefore never attaches an ERC purely because it
happens to share species with the current module and produces no
productive novelty: Mode 1 only attaches producers of an \emph{open}
requirement, and Mode 2's three moves are all instances of productive
novelty (synergy, complementarity, or absorbing a descendant already
present). The result is that Algorithm~\ref{alg:main} computes exactly
$\SOgen$, not $\SO$ -- which is the escape from the combinatorial
explosion, since the latent joins of otherwise-unrelated modules already
number exponentially in a network with several independent persistent
components. A latent module is reconstructed on demand in $O(1)$ closure
computations rather than ever being searched for:

\begin{algorithm}[t]
\caption{LatentJoin($X$, $Y$)}\label{alg:latent}
\begin{algorithmic}[1]
\State $Z \leftarrow \cl(X\,|\,Y)$
\If{$\req(Z)=\emptyset$ and $Z$ is connected}
  \State \Return $Z$
\Else
  \State \Return \textsc{None}
\EndIf
\end{algorithmic}
\end{algorithm}

\paragraph{Complexity.} Each Mode-1/Mode-2 state transition is $O(1)$
amortised bitmask work per ERC touched by Algorithm~\ref{alg:erc-syn-close}
(no reaction scan). Branching per open state is bounded by the number of
fundamental synergy partners and $\subseteq$-minimal producers/consumers
touching the state -- the fundamental-relation density -- never by $|\E|$
or $\binom{|\E|}{2}$; the companion study's empirical finding that this
density \emph{shrinks} as networks grow~\cite{BMBcompanion} is what makes
the search tractable at genome scale (Section~\ref{sec:experiments}). The
one remaining global-complexity step is order assignment
(line 8 of Algorithm~\ref{alg:main}), a containment sweep over all
discovered SOs found so far: $O(|\SOgen|^2)$ in the worst case, the
dominant cost on networks with thousands of ESPMs
(Section~\ref{sec:experiments}).

% =============================================================================
\section{Bitset Implementation}\label{sec:implementation}
% =============================================================================

All species and ERC sets are represented as Python \texttt{int} values
used as $n$-bit integers (species $s_i\leftrightarrow$ bit $i$), giving
union, intersection, containment, and complement as single machine
bitwise operations; Python integers are arbitrary precision, so networks
with more than 64 species are handled transparently.

\begin{itemize}
  \item \textbf{ERCData} (frozen, one per ERC): \texttt{species\_mask},
  \texttt{req\_mask}, \texttt{prod\_mask}, \texttt{min\_bases}
  (list of bitsets), \texttt{reaction\_indices}.
  \item \textbf{HierarchyData}: per-ERC \texttt{parents} (direct Hasse
  covers), \texttt{children}, \texttt{ancestors}, \texttt{descendants}
  (transitively closed frozensets).
  \item \textbf{FundamentalGraph}: precomputed once per network from the
  three structures above --
  \texttt{syn\_from}$[i]=\{(j,k):E_i+E_j\twoheadrightarrow E_k\}$
  (stored symmetrically in $i$ and $j$),
  \texttt{comp\_by\_species}$[s]=\minprod(s)$,
  \texttt{comp\_consumers\_by\_species}$[s]=\mincons(s)$.
  \item \textbf{DFSState} (immutable): $(\mathrm{erc\_set},\ \mathrm{sp},\
  \mathrm{req},\ \mathrm{prod},\ \mathrm{min\_ext})$ -- everything
  Algorithm~\ref{alg:erc-syn-close} needs, and everything the search
  visits or dedups on.
\end{itemize}

The whole EPM/ESPM search touches the reaction network only twice: once
to build $\E$ (Section~\ref{sec:stage1}) and once, for each ERC, to
precompute $\req$/$\prods$ from $\R_E$. Every subsequent state transition
is pure ERC-index-space bitmask arithmetic.

% =============================================================================
\section{Oracle-First Validation}\label{sec:oracle}
% =============================================================================

Each stage has an independent brute-force oracle used exclusively for
cross-validation, never in the production pipeline:

\begin{itemize}
  \item \texttt{erc\_oracle}: enumerates all $X\subseteq\M$ and checks
  closure directly -- exponential in $|\M|$, usable for small hand-built
  networks.
  \item \texttt{synergy\_oracle}: a plain triple loop, filtered
  basic$\to$maximal$\to$fundamental exactly per
  Definitions~\ref{def:fund_syn} (used to validate both
  Algorithm~\ref{alg:syn} and Algorithm~\ref{alg:syn-basis}).
  \item \texttt{comp\_basic\_set} / \texttt{comp\_fund\_set}: plain pair
  and species loops computing basic complementarity and
  $\minprod(s)$/$\mincons(s)$ directly from Definitions~3--5.
  \item \texttt{epm\_oracle} / \texttt{espm\_oracle}: enumerate all
  $2^{|\E|}$ ERC subsets, compute the joint closure of each, and keep
  those that are closed, SSM, reactive, and connected; the $\subseteq$-minimal
  such sets are the EPMs, and order is assigned by the same containment
  rule as Definition~\ref{def:order}. Feasible up to $|\E|\lesssim 24$.
\end{itemize}

\subsection{Validation methodology and results}

The EPM/ESPM search was validated by exact-set comparison against
\texttt{epm\_oracle}/\texttt{espm\_oracle} on (i) five hand-built networks
covering loops, hierarchies, non-persistence, and inflow quotienting;
(ii) a fully worked network (Section~\ref{sec:worked}) exhibiting every
mechanism, including a sterile ERC; and (iii) $25$ real BioModels
networks with $4$--$24$ ERCs, the size at which the oracle itself remains
tractable. Because Algorithm~\ref{alg:main} deliberately searches only
$\SOgen$ (Section~\ref{sec:grammar}), agreement with the oracle -- which
enumerates all of $\SO$ -- is defined modulo latent joins: every module
found is required to be a genuine member of $\SO$ (soundness), and every
oracle module absent from the result is required to be reconstructible by
repeatedly applying Algorithm~\ref{alg:latent} to modules already found
(completeness of the generated core). Both conditions held exactly on all
$30$ test networks. This process caught two real defects during
development, both now fixed and covered by regression tests: (a) an
earlier version of Mode 2 seeded growth past an SSM only through synergy
edges, so it silently returned zero ESPMs on any network whose novelty was
purely complementarity- or lift-driven -- exactly the worked example of
Section~\ref{sec:worked}; and (b) fixing (a) exposed
Lemma~\ref{lem:synignite}'s membership-vs-subset distinction, an earlier
version of Algorithm~\ref{alg:erc-syn-close} tested ERC-index membership
rather than species-subset containment and briefly produced states whose
tracked masks disagreed with their true reaction-level closure on a real
BioModels network before the fix.

\subsection{A fully worked network}
\label{sec:worked}

Let $\M=\{a,b,c,d,e,f,g,h,k,p,q,w,z\}$ and
\begin{align*}
r_A&=f{+}a\to 2a{+}p, & r_B&=p{+}b\to 2b{+}f, & r_C&=p{+}c\to 2c{+}f,\\
r_D&=g{+}d\to 2d{+}q, & r_E&=q{+}e\to 2e{+}g, & r_F&=p{+}h\to 2h,\\
r_H&=p{+}z\to z{+}b, & r_G&=k{+}w\to 2k{+}p. &&
\end{align*}
This gives eight ERCs (Table~\ref{tab:worked}); the only containment is
$E'\supset E_B$ ($E'=\cl\{p,z\}=\{p,z,b,f\}$, since $r_H$ makes $b$, which
ignites $r_B$). There are no synergies -- every reaction ignites from a
single ERC's closure -- so every relation here is complementary, the
cleanest setting to exercise Mode 2.

\begin{table}[t]
\centering\small
\begin{tabular}{lll}
\toprule
ERC & species & $\req$ \\
\midrule
$E_A=\cl\{f,a\}$ & $\{f,a,p\}$ & $\{f\}$\\
$E_B=\cl\{p,b\}$ & $\{p,b,f\}$ & $\{p\}$\\
$E_C=\cl\{p,c\}$ & $\{p,c,f\}$ & $\{p\}$\\
$E_D=\cl\{g,d\}$ & $\{g,d,q\}$ & $\{g\}$\\
$E_E=\cl\{q,e\}$ & $\{q,e,g\}$ & $\{q\}$\\
$E^{*}=\cl\{p,h\}$ & $\{p,h\}$ & $\{p\}$\\
$E'=\cl\{p,z\}$ & $\{p,z,b,f\}$ & $\{p\}$\\
$E_{\dagger}=\cl\{k,w\}$ & $\{k,w,p\}$ & $\{w\}$\\
\bottomrule
\end{tabular}
\caption{ERCs of the worked network. $E_\dagger$ produces $p$ but requires
$w$, which no reaction produces: Mode 1 seeded from $E_\dagger$ dead-ends
the moment $w$ has no complementary producer, so $E_\dagger$'s species
never appear in any discovered module -- with no separate sterility
fixpoint required to establish this in advance.}
\label{tab:worked}
\end{table}

Algorithm~\ref{alg:main} finds exactly the three EPMs
$Y_1=E_A\js E_B=\{f,a,p,b\}$, $Y_2=E_A\js E_C=\{f,a,p,c\}$, and (a disjoint
component) $Y_3=E_D\js E_E=\{g,d,q,e\}$, and exactly the following eight
ESPMs, matching \texttt{espm\_oracle} exactly:
\begin{align*}
\text{order 1:}\quad
&\{f,a,p,b,c\}\ (=Y_1\js Y_2,\text{ a latent join}),\\
&\{f,a,p,b,h\}\ (=Y_1\js E^{*},\text{ consumer complementarity}),\\
&\{f,a,p,c,h\}\ (=Y_2\js E^{*}),\\
&\{f,a,p,b,z\}\ (=Y_1\js E',\text{ vertical lift over }E_B);\\
\text{order 2:}\quad
&\{f,a,p,b,c,h\},\ \{f,a,p,b,c,z\},\ \{f,a,p,b,h,z\};\\
\text{order 3:}\quad
&\{f,a,p,b,c,h,z\}.
\end{align*}
$Y_1\js Y_2$ (order 1) is discovered as a plain consumer-complementarity
attachment: $E_C$ is a fundamental consumer of $p$, which $Y_1$ already
produces via $E_A$, so it is found by exactly the same mechanism as the
genuine discoveries -- Definition~9's ``latent'' classification is a
statement about \emph{whether} a step was productively novel along some
irreducible generator, not about which search mechanism reaches a module,
and $Y_1\js Y_2$ is simultaneously reachable as a latent join and as a
consumer-complementarity attachment; excluding it from the search would
require tracking generator irreducibility explicitly, which the current
implementation does not do (Section~\ref{sec:conclusion}). $\{f,a,p,b,z\}$
is the module that specifically requires the vertical-lift move
(item~2 of Section~\ref{subsec:moves}): $E'$'s minimal producer of $p$ is
$E_A$, already present, but $E'$ itself is not a \emph{minimal} consumer
of $p$ (its sub-ERC $E_B$ is, and is already present too), so
Algorithm~\ref{alg:fund_comp} never registers $E'$ as a fundamental
complementarity partner of anything -- the only route to it is lifting
directly from its hierarchy parent relation to $E_B$.

% =============================================================================
\section{Experiments}\label{sec:experiments}
% =============================================================================

We report one representative genome-scale run rather than a full-corpus
sweep (left for future work, Section~\ref{sec:conclusion}). On the BiGG
reconstruction \texttt{iNJ661} (825 species, 1296 reactions), the pipeline
finds $712$ ERCs ($210$ persistent) in $3.7\,\mathrm{s}$; the containment
hierarchy in $0.06\,\mathrm{s}$; $3{,}697$ fundamental synergies (of
$240{,}833$ basic, $123{,}451$ maximal) via Algorithm~\ref{alg:syn-basis}
in $2.9\,\mathrm{s}$; $1{,}097$ fundamental complementarities (of $6{,}738$
basic, $3{,}139$ pure) in $0.2\,\mathrm{s}$; $329$ primitive ERCs reaching
$87.9\%$ synergy coverage of $\E$ in $1\,\mathrm{ms}$; and
Algorithm~\ref{alg:main} finds all $72$ EPMs and $1{,}124$ ESPMs across
$17$ orders in $12.5\,\mathrm{s}$ ($17{,}997$ states explored: $1{,}124$
SSMs, $8{,}271$ dead ends, the remainder branching states), for a total
pipeline time of $19.4\,\mathrm{s}$ on a single desktop core -- a regime
in which the $2^{712}$-subset brute-force oracle is entirely infeasible.

% =============================================================================
\section{Conclusion}\label{sec:conclusion}
% =============================================================================

We have described, implemented, and validated a complete pipeline for the
generative structure of chemical organizations: ERC discovery, the
containment hierarchy, basic/maximal/fundamental synergy (with a
basis-first algorithm markedly faster than pair-first on sparse
networks), basic/pure/fundamental complementarity, primitive-ERC synergy
reachability, and -- the paper's central contribution -- exact
computation of every EPM and every ESPM of any order, via a two-mode
ERC-index-space traversal that never recomputes a reaction-level closure
and deliberately confines its search to the generated core, excluding
latent joins in favour of an $O(1)$ on-demand reconstruction. This closes
the gap left open by the companion existence theory~\cite{BMBcompanion}
and by an earlier version of this pipeline that covered only single-ERC
and pairwise EPMs; correctness is established not by proof sketch alone
but by exhaustive oracle cross-validation on hand-built, worked, and real
networks.

Concrete extensions this implementation makes natural, and which we leave
as future work: (i) precomputing the requirement-reachability fixpoint
that identifies \emph{sterile} ERCs -- those whose requirement can never
be reduced to $\emptyset$ by any reachable combination -- so that
provably dead branches are pruned before rather than after being
explored, as sketched (but not yet implemented) in the companion
algorithmic theory; (ii) an explicit generator-irreducibility check so
that latent joins reachable through consumer complementarity (as in
Section~\ref{sec:worked}) are recognised and excluded from the search
rather than incidentally found by it; (iii) grade tracking -- the lowest
order at which an ERC first becomes load-bearing -- as a structural,
per-module summary statistic; (iv) replacing the $O(|\SOgen|^2)$ post-hoc
order-assignment sweep with an incremental, order-indexed dedup structure
for networks with tens of thousands of ESPMs; and (v) a full-corpus
benchmark, in the style of the companion study's $438$-network survey,
now measuring computation rather than counts.

\begin{thebibliography}{9}
\bibitem{BMBcompanion}
T.~Veloz and A.~Bassi,
``Synergy and Complementarity: The Generative Basis of Chemical
Organizations,''
\textit{Bulletin of Mathematical Biology}, in press.

\bibitem{dittrich2007chemical}
P.~Dittrich and P.~Speroni~di~Fenizio,
``Chemical Organisation Theory,''
\textit{Bulletin of Mathematical Biology}, vol.~69, no.~4,
pp.~1199--1231, 2007.

\bibitem{dowling1984linear}
W.~F.~Dowling and J.~H.~Gallier,
``Linear-time algorithms for testing the satisfiability of propositional
Horn formulae,''
\textit{Journal of Logic Programming}, vol.~1, no.~3, pp.~267--284, 1984.

\bibitem{bigg2019}
Z.~King \textit{et~al.},
``BiGG Models: A platform for integrating, standardizing, and sharing
genome-scale models,''
\textit{Nucleic Acids Research}, vol.~44, no.~D1, pp.~D515--D522, 2016.
\end{thebibliography}

\end{document}
